% VI27-DETAILED-PROBLEM-09
\begin{problem}[Regular locus of a rational function]\label{prob:vi27-p09}
For \(r\in K(X)\), prove that the points where \(r\) is regular form a nonempty open subset.
\par\smallskip\noindent\textit{Provenance.} Canonical VI/27 dossier.
\end{problem}

\ifdefined\FullProblemDossiers
\begin{solution}
Choose nonempty affine \(U=\Spec A\) and write \(r=a/b\in\Frac(A)\) with \(b\ne0\). On \(D(b)\), \(b\) is
invertible, so \(r\in A_b=\mathcal O_X(D(b))\). Thus the regular locus is nonempty. If \(r\) is regular at a
point \(x\), then by the stalk definition it is represented by a section on some neighborhood \(W_x\) of \(x\);
the same section makes \(r\) regular at every point of \(W_x\). Hence the regular locus is the union of the
\(W_x\) and is open.
\end{solution}
\fi
